\documentclass[12pt]{article}
\usepackage[utf8]{inputenc}
\usepackage[T1]{fontenc}
\usepackage{hyperref}
\title{Kwaza Language Resources for K-12 Curriculum Development}
\author{Compiled by Claude Sonnet 5.5 (claude-sonnet-5-5) for Andr\'as Kornai}
\date{}
\begin{document}
\maketitle

\section{Introduction}
Kwaza (ISO 639-3 xwa; Glottocode kwaz1243) is a language isolate of southern Rond\^onia, Brazil, with very few speakers. The count is about 7 in the task description, 25 in the publisher's blurb and 54 in a 2004 count cited on Wikipedia. Earlier documentation was limited to short word lists. The main documentation is Hein van der Voort's grammar.

\section{Primary Educational Resources}
No primers, textbooks or school programmes were found in the searches.

\section{Linguistic and Reference Materials}
\begin{itemize}
\item van der Voort, \emph{A Grammar of Kwaza} (Mouton de Gruyter, 2004). It includes texts and a Kwaza dictionary with an English--Kwaza finder list: \url{https://www.degruyterbrill.com/document/doi/10.1515/9783110197280/html}
\item Catalogue record of the same work: \url{https://oeaw.ac.at/resources/Record/9783110197280}
\item Wikipedia, Kwaza language: \url{https://en.wikipedia.org/wiki/Kwaza_language}
\end{itemize}

\section{Religious and Community Resources}
None found.

\section{Digital and Audio Resources}
The grammar was published with an audio CD, according to a bookseller listing. No online audio was found.

\section{Curriculum Development Context}
Kwaza is a tiny, highly endangered language. Existing materials are scholarly, so any school programme would draw on the grammar and texts, with community involvement. No ongoing school use was found.
\end{document}